%! Potential Errors When Performing Tests — Unit 6, Lesson 6.7 — Homework
\documentclass[10pt]{article}
\usepackage{apstats-article}
\usepackage{apstats-boxes}
\usepackage{grading}
\gradingassignment{apstats-u06-l07-hw}
\gradingcourse{advanced_statistics}
\newcommand{\TallMath}[1]{$\displaystyle #1\rule[-1.4em]{0pt}{3.2em}$}

\begin{document}

\pageheader{Unit 6, Lesson 6.7}{Homework}

\vspace{0.15in}

\begin{enumerate}

  \item A city health department tests $H_0: p = 0.10$ (10\% of adults have high blood
        pressure) vs.\ $H_a: p > 0.10$.
        \begin{enumerate}[label=\alph*.]
          \item Describe a Type I error in context. \gradedhide{\writeline\\\gradedhere{\gradedblank[id=hw1a, type=rubric, points=2, display={Concluding more than 10\% of the city's adults have high blood pressure when the true rate is 10\%.}]{\linewidth+6pt}{\gradedlinesheight{1}}}\writelines{1}}
          \item Describe a Type II error in context. \gradedhide{\writeline\\\gradedhere{\gradedblank[id=hw1b, type=rubric, points=2, display={Failing to conclude more than 10\% have high blood pressure when the true rate really exceeds 10\%.}]{\linewidth+6pt}{\gradedlinesheight{1}}}\writelines{1}}
          \item If the true proportion is 12\%, would power be higher or lower than if
                the true proportion were 20\%? Explain. \gradedhide{\writeline\\\gradedhere{\gradedblank[id=hw1c, type=rubric, points=2, display={Lower at 12\%: 0.12 is closer to $p_0 = 0.10$ than 0.20 is, so the difference is harder to detect.}]{\linewidth+6pt}{\gradedlinesheight{1}}}\writelines{1}}
        \end{enumerate}

  \item A new reading intervention is tested at a school. $H_0$: the intervention does not
        improve reading scores; $H_a$: the intervention improves reading scores.
        \begin{enumerate}[label=\alph*.]
          \item From the perspective of a school principal, which error would be more
                costly? Explain. \gradedhide{\writeline\\\gradedhere{\gradedblank[id=hw2a, type=rubric, points=2, display={Type II: an intervention that really works would not be adopted, so students lose a real benefit (a Type I error only wastes resources).}]{\linewidth+6pt}{\gradedlinesheight{2}}}\writelines{2}}
          \item The school increases the sample from 60 students to 240 students. How
                does this change the power of the test and Type II error probability? \gradedhide{\writeline\\\gradedhere{\gradedblank[id=hw2b, type=rubric, points=2, display={Power increases and $P(\text{Type II error})$ decreases (larger $n$, smaller $SE$).}]{\linewidth+6pt}{\gradedlinesheight{1}}}\writelines{1}}
        \end{enumerate}

  \item \textbf{True or False.} For each statement, write T or F and explain.
        \begin{enumerate}[label=\alph*.]
          \item ``Reducing $\alpha$ from 0.05 to 0.01 always increases power.''
                \rlap{\smash{\gradedblank[id=hw3a1, type=choice, points=1, answer=F]{1cm}{0.17in}}}\phantom{\blank{1cm}} Explain: \gradedhide{\writeline\\\gradedhere{\gradedblank[id=hw3a2, type=rubric, points=1, display={False: a smaller $\alpha$ makes $H_0$ harder to reject, so power decreases.}]{\linewidth+6pt}{\gradedlinesheight{1}}}\writelines{1}}
          \item ``A Type II error can only occur if $H_0$ is rejected.''
                \rlap{\smash{\gradedblank[id=hw3b1, type=choice, points=1, answer=F]{1cm}{0.17in}}}\phantom{\blank{1cm}} Explain: \gradedhide{\writeline\\\gradedhere{\gradedblank[id=hw3b2, type=rubric, points=1, display={False: a Type II error happens when we FAIL to reject $H_0$ while $H_a$ is true.}]{\linewidth+6pt}{\gradedlinesheight{1}}}\writelines{1}}
        \end{enumerate}

\end{enumerate}

\begin{reflectionbox}
\textbf{Preview:} Lesson 6.8 introduces the two-sample $z$-interval for the difference
of two proportions $p_1 - p_2$ --- when we want to compare two groups rather than test
a single proportion against a benchmark.
\end{reflectionbox}

\end{document}
